\documentclass[10pt,a4paper]{amsart}
\usepackage[margin=19mm]{geometry}
\usepackage{amsmath,amssymb,mathtools}
\usepackage[scr=rsfs,cal=euler]{mathalfa}
\usepackage{xfrac}
\usepackage[unicode,hidelinks,bookmarks=false]{hyperref}
\newcommand{\ie}{i.e.\ }
\newcommand{\cf}{cf.\ }
\newcommand{\resp}{respectively\ }
\newcommand{\Subsection}{\subsection}
\providecommand{\nobreakdash}{\nobreak\hbox{-}}
\setlength{\emergencystretch}{2em}
\multlinegap0pt
\usepackage{amssymb}
\usepackage{tikz}
\usepackage{mdframed}
\usepackage{mathtools}
\usepackage{bbm}
\usepackage[shortlabels]{enumitem}
\usepackage[normalem]{ulem}
\usepackage{xcolor}
\newtheorem{proposition}{Proposition}
\title{Rates of Bulk Convergence for Ensembles of Classical Compact Groups}
\author{Mengchun Cai}
\date{Source: arXiv:2508.21274v2 (CC BY 4.0)}
\begin{document}
\maketitle

\noindent\emph{Source and licence.} Rates of Bulk Convergence for Ensembles of Classical Compact Groups. Version arXiv:2508.21274v2 (CC BY 4.0), distributed by arXiv under the Creative Commons Attribution 4.0 International licence, http://creativecommons.org/licenses/by/4.0/, as recorded in the arXiv licence metadata for this exact version and shown on the abstract page https://arxiv.org/abs/2508.21274v2. Full licence notice: \texttt{licenses/arxiv-2508.21274-CC-BY-4.0.txt}.\par\smallskip

\noindent\textit{Editorial scope.} Counted unit: the article's proposition fact (source0.tex lines 235--266). The article's complete original proof of it is delivered with it (source0.tex lines 267--290, one \texttt{proof} environment of 24 lines). No mathematics is written, repaired or re-derived: the statement and the proof are the article's own lines, in the article's own notation, and no retained line is substituted or re-flowed.  No further source-local context is needed: every cross-reference of every retained line resolves inside the two delivered spans.  Nothing was omitted from any retained span, so the package carries no omission record.  No retained line contains a citation, so the wrapper carries no bibliography and no bibliography entry.  No retained line contains a figure environment, an image-inclusion command or drawing code, so no image is delivered and the package declares no asset dependency.  Editorial formatting only, in addition: an AMS \texttt{amsart} preamble that restates the article's own declarations and macros used by the retained lines, namely the environments \texttt{proposition} on separate counters, together with the standard packages those lines need.  The retained environments print in this document as Proposition 1 in order of appearance, whereas the article prints the same items with the numbering of its own full document.  The complete native source of the article as distributed in the arXiv e-print for this exact version is bundled unchanged as \texttt{sources/184099/source0.tex}.
\begin{proposition}\label{CUE fact}
Let $I=[-s,s]$ for some $s>0$, $\mathcal{A}_{2k+1}$ and $\mathcal{A}'_{2k+1}$ be operators on $L^2(I)$ with kernels:
\begin{equation*}
A_{2k+1}(x,y)=(\pi(x-y))^{2k+1}\sin(\pi(x-y))
\end{equation*}
and
\begin{equation*}
A'_{2k+1}(x,y)=(\pi(x+y))^{2k+1}\sin(\pi(x+y))
\end{equation*}
respectively.
We have
\begin{equation*}
\mathcal{A}_{2k+1}=\sum_{j=0}^{2k+1}\left(\begin{matrix}
    2k+1\\
    j
\end{matrix}\right)M_{S^+_{j}}\mathcal{C}^-_{2k+1-j}-\sum_{j=0}^{2k+1}\left(\begin{matrix}
    2k+1\\
    j
\end{matrix}\right)M_{C^+_{j}}\mathcal{S}^-_{2k+1-j}
\end{equation*}
and
\begin{equation*}
\mathcal{A}'_{2k+1}=\sum_{j=0}^{2k+1}\left(\begin{matrix}
    2k+1\\
    j
\end{matrix}\right)M_{S^+_{j}}\mathcal{C}^+_{2k+1-j}+\sum_{j=0}^{2k+1}\left(\begin{matrix}
    2k+1\\
    j
\end{matrix}\right)M_{C^+_{j}}\mathcal{S}^+_{2k+1-j}
\end{equation*}
where $M_f:g\rightarrow f\cdot g$ is the multiplication operator.
\end{proposition}

\begin{proof}
Directly applying the binomial formula and trigonometric identities to $A_{2k+1}$ leads to
\begin{equation}\label{S4A}
    \begin{split}
        &A_{2k+1}(x,y)=\sum_{j=0}^{2k+1}\left(\begin{matrix}
            2k+1\\j
        \end{matrix}\right)(\pi x)^j(-\pi y)^{2k+1-j}\sin(\pi x)\cos(\pi y)\\&~~~~~~~~~~~~~-\sum_{j=0}^{2k+1}\left(\begin{matrix}
            2k+1\\j
        \end{matrix}\right)(\pi x)^j(-\pi y)^{2k+1-j}\sin(\pi y)\cos(\pi x)\\
        &=\sum_{j=0}^{2k+1}\left(\begin{matrix}
            2k+1\\j
        \end{matrix}\right)S_{j}^+(x)C_{2k+1-j}^{-}(y)-\sum_{j=0}^{2k+1}\left(\begin{matrix}
            2k+1\\j
        \end{matrix}\right)C_{j}^+(x)S_{2k+1-j}^-(y).
    \end{split}
\end{equation}
For each $j$, assume $\tilde{K}_j(x,y)=S_j^+(x)C_{2k+1-j}^-(y)$ is the integral kernel of $\tilde{\mathcal{K}}_{j}$ on $L^2(I)$. Then for any $f\in L^2(I)$, 
\begin{equation*}
    \tilde{\mathcal{K}}_jf(x)=\int_I S_j^+(x)C_{2k+1-j}^-(y)f(y)dy=M_{S_j^+}\left(\int_I C_{2k+1-j}^-(y)f(y)dy\right)=M_{S_j^+}\mathcal{C}_{2k+1-j}^-f(x)
\end{equation*}
and hence, $\tilde{\mathcal{K}}_j=M_{S_j^+}\mathcal{C}_{2k+1-j}^-$. Similar discussion for $C^+_j(x)S_{2k+1-j}^-(y)$, together with the linearity of the integral, yields the decomposition for $\mathcal{A}_{2k+1}$ from \eqref{S4A}.

Substituting all negative signs by the positive ones in the previous discussion for $A_{2k+1}(x,y)$ completes the proof for $\mathcal{A}'_{2k+1}$.
\end{proof}

\end{document}
